\sgaStableTarget{sga6.e09.lem.1-7}\textbf{Lemma 1.7.} \uline{Let} @@M0@@ \uline{be a principal ideal domain}, @@M1@@. \uline{Then the canonical homomorphism}:
@@M2@@
\uline{is an isomorphism}.

Since @@M3@@ is regular, the homomorphism @@M4@@ is an isomorphism. Thus @@M5@@ is generated by the classes of locally free @@M6@@-Modules of finite type, i.e.\ by the classes of projective @@M7@@-modules of finite type. Since @@M8@@ is a principal ideal domain, every projective module of finite type is free, giving surjectivity. Injectivity follows from the fact that @@M9@@ is @@M10@@-augmented.

\sgaStableTarget{sga6.e09.prop.1-8}\textbf{Proposition 1.8.} \uline{Let} @@M11@@ \uline{be a noetherian scheme}, @@M12@@ \uline{a locally free} @@M13@@\uline{-Module of rank} @@M14@@, @@M15@@ \uline{the vector bundle associated to} @@M16@@, @@M17@@ \uline{the structural morphism of} @@M18@@, \uline{and} @@M19@@ \uline{the zero section of} @@M20@@. \uline{Then the composite homomorphism}:
@@M21@@
(\uline{where} @@M22@@ \uline{is defined by} IV 2.12.1 \uline{and} 1.5) \uline{is multiplication by the element} @@M23@@ \uline{of} @@M24@@, \uline{denoting by} @@M25@@ \uline{the class of} @@M26@@ \uline{in} @@M27@@, \uline{and putting}
@@M28@@

This statement is the analogue, for @@M29@@, of VII 2.7.

Let @@M30@@ be a coherent @@M31@@-Module, of class @@M32@@ in @@M33@@. Then @@M34@@ is the class in @@M35@@ of the complex @@M36@@, hence of @@M37@@ regarded as an @@M38@@-Module by means of @@M39@@. To compute @@M40@@, one must therefore take the class of @@M41@@.
